\FigureLayoutDeclare{img/2006/zhejiang_liberal/q16_fig1.png}{7.0cm}{69bp 43bp 33bp 29bp}
\FigureLayoutDeclare{img/2006/zhejiang_liberal/q19_fig1.png}{9.0cm}{27bp 38bp 32bp 34bp}

\chapter{2006年浙江卷（文）}
\section{单选题}
\begin{problem}
设集合 \(A = \{x \mid -1 \leqslant x \leqslant 2\}\)，\(B = \{x \mid 0 \leqslant x \leqslant 4\}\)，则 \(A \cap B =\)（\quad）

\choices
    {\([0,2]\)}
    {\([1,2]\)}
    {\([0,4]\)}
    {\([1,4]\)}
\end{problem}

\begin{answer}
A
\end{answer}

\begin{solution}
由交集的定义，\(x\) 必须同时满足 \(-1\leqslant x\leqslant2\) 和 \(0\leqslant x\leqslant4\)，所以 \(0\leqslant x\leqslant2\)，即 \(A\cap B=[0,2]\)，故选 A.
\end{solution}

\begin{problem}
在二项式 \((x + 1)^6\) 的展开式中，含 \(x^3\) 的项的系数是（\quad）

\choices
    {\(15\)}
    {\(20\)}
    {\(30\)}
    {\(40\)}
\end{problem}

\begin{answer}
B
\end{answer}

\begin{solution}
在展开式中，含 \(x^3\) 的项是从六个因式中取三个 \(x\) 和三个 \(1\) 所得到的项，其系数为 \(\mathrm C_6^3=20\)，故选 B.
\end{solution}

\begin{problem}
抛物线 \( y^{2} = 8x \) 的准线方程是（\quad）

\choices
    {\(x = -2\)}
    {\(x = -4\)}
    {\(y = -2\)}
    {\(y = -4\)}
\end{problem}

\begin{answer}
A
\end{answer}

\begin{solution}
抛物线的标准形式为 \(y^2=2px\)，与 \(y^2=8x\) 比较得 \(2p=8\)，所以 \(p=4\). 抛物线 \(y^2=2px\) 的准线为 \(x=-\frac p2\)，故准线方程为 \(x=-2\)，选 A.
\end{solution}

\begin{problem}
已知 \(\log_{\frac{1}{2}} m < \log_{\frac{1}{2}} n < 0\)，则（\quad）

\choices
    {\(n < m < 1\)}
    {\(m < n < 1\)}
    {\(1 < m < n\)}
    {\(1 < n < m\)}
\end{problem}

\begin{answer}
D
\end{answer}

\begin{solution}
因为 \(0<\frac12<1\)，函数 \(y=\log_{\frac12}x\) 在 \((0,+\infty)\) 上单调递减. 由 \(\log_{\frac12}m<\log_{\frac12}n\)，得 \(m>n\). 又由 \(\log_{\frac12}n<0=\log_{\frac12}1\)，得 \(n>1\). 因而 \(1<n<m\)，故选 D.
\end{solution}

\begin{problem}
设向量 \(\bs{a}\)，\(\bs{b}\)，\(\bs{c}\) 满足 \(\bs{a} + \bs{b} + \bs{c} = 0\)，\(\bs{a} \perp \bs{b}\)，\(|\bs{a}| = 1\)，\(|\bs{b}| = 2\)，则 \(|\bs{c}|^2 =\)（\quad）

\choices
    {\(1\)}
    {\(2\)}
    {\(4\)}
    {\(5\)}
\end{problem}

\begin{answer}
D
\end{answer}

\begin{solution}
由 \(\bs a+\bs b+\bs c=0\)，得 \(\bs c=-(\bs a+\bs b)\). 因为 \(\bs a\perp\bs b\)，所以 \(\lvert\bs c\rvert^2=\lvert\bs a+\bs b\rvert^2=\lvert\bs a\rvert^2+\lvert\bs b\rvert^2+2\bs a\cdot\bs b=1^2+2^2=5\)，故选 D.
\end{solution}

\begin{problem}
函数 \( f(x) = x^3 - 3x^2 + 2 \) 在区间 \([-1, 1]\) 上的最大值是（\quad）

\choices
    {\(-2\)}
    {\(0\)}
    {\(2\)}
    {\(4\)}
\end{problem}

\begin{answer}
C
\end{answer}

\begin{solution}
求导得 \(f'(x)=3x^2-6x=3x(x-2)\). 在 \([-1,0]\) 上 \(f'(x)\geqslant0\)，在 \([0,1]\) 上 \(f'(x)\leqslant0\)，所以最大值只需比较 \(x=-1,0,1\) 处的函数值. 计算得 \(f(-1)=-2\)，\(f(0)=2\)，\(f(1)=0\)，故最大值为 \(2\)，选 C.
\end{solution}

\begin{problem}
“\(a > 0\)，\(b > 0\)”是“\(ab > 0\)”的（\quad）

\choices
    {充分而不必要条件}
    {必要而不充分条件}
    {充分必要条件}
    {既不充分也不必要条件}
\end{problem}

\begin{answer}
A
\end{answer}

\begin{solution}
若 \(a>0\)，\(b>0\)，则 \(ab>0\)，所以“\(a>0\)，\(b>0\)”是“\(ab>0\)”的充分条件. 但它不是必要条件，例如 \(a=-1\)，\(b=-1\) 时 \(ab=1>0\)，而 \(a>0\)，\(b>0\) 不成立，故选 A.
\end{solution}

\begin{problem}
如图，正三棱柱 \(ABC - A_{1}B_{1}C_{1}\) 的各棱长都 \(2\)，\(E\)，\(F\) 分别是 \(AB\)，\(A_{1}C_{1}\) 的中点，则 \(EF\) 的长是（\quad）

\bitmapfigure[max width=0.42\linewidth]{img/2006/zhejiang_liberal/q8_fig1.png}

\choices
    {\(2\)}
    {\(\sqrt{3}\)}
    {\(\sqrt{5}\)}
    {\(\sqrt{7}\)}
\end{problem}

\begin{answer}
C
\end{answer}

\begin{solution}
建立直角坐标系，令 \(A=(0,0,0)\)，\(B=(2,0,0)\)，\(C=(1,\sqrt3,0)\)，并令 \(A_1=(0,0,2)\)，\(C_1=(1,\sqrt3,2)\). 则 \(E=(1,0,0)\)，\(F=\left(\frac12,\frac{\sqrt3}{2},2\right)\)，从而 \(\lvert EF\rvert^2=\left(-\frac12\right)^2+\left(\frac{\sqrt3}{2}\right)^2+2^2=5\). 因此 \(\lvert EF\rvert=\sqrt5\)，选 C.
\end{solution}

\begin{problem}
在平面直角坐标系中，不等式组 \(\begin{cases}
x+y-2\geqslant0,\\
x-y+2\geqslant0,\\
x\leqslant2
\end{cases}\) 表示的平面区域的面积是（\quad）

\choices
    {\(4\sqrt{2}\)}
    {\(4\)}
    {\(2\sqrt{2}\)}
    {\(2\)}
\end{problem}

\begin{answer}
B
\end{answer}

\begin{solution}
不等式组等价于 \(y\geqslant2-x\)、\(y\leqslant x+2\) 和 \(x\leqslant2\). 两条边界直线的交点为 \((0,2)\)，直线 \(x=2\) 与它们的交点为 \((2,0)\) 和 \((2,4)\). 因而区域是以这三点为顶点的三角形，其底边长为 \(4\)，到底边的高为 \(2\)，面积为 \(\frac12\times4\times2=4\)，故选 B.
\end{solution}

\begin{problem}
对 \(a\)，\(b \in \R\)，记 \(\max \{a, b\} = \begin{cases}
a, & a \geqslant b,\\
b, & a < b,
\end{cases}\) 函数 \(f(x) = \max \{|x + 1|, |x - 2|\}\) （\(x \in \R\)）的最小值是（\quad）

\choices
    {\(0\)}
    {\(\frac{1}{2}\)}
    {\(\frac{3}{2}\)}
    {\(3\)}
\end{problem}

\begin{answer}
C
\end{answer}

\begin{solution}
对任意实数 \(x\)，由三角不等式得 \(\lvert x+1\rvert+\lvert x-2\rvert\geqslant3\). 因此 \(\max\{\lvert x+1\rvert,\lvert x-2\rvert\}\geqslant\frac{\lvert x+1\rvert+\lvert x-2\rvert}{2}\geqslant\frac32\). 当 \(x=\frac12\) 时，\(\lvert x+1\rvert=\lvert x-2\rvert=\frac32\)，可以取得上述下界. 所以最小值为 \(\frac32\)，选 C.
\end{solution}

\section{填空题}
\begin{problem}
不等式 \(\frac{x + 1}{x - 2} > 0\) 的解集是\fillinblank{}.
\end{problem}

\begin{answer}
\(\{x \mid x < -1\text{ 或 }x > 2\}\)
\end{answer}

\begin{solution}
临界点为 \(x=-1\) 和 \(x=2\)，其中 \(x=2\) 不在定义域内. 在 \((-\infty,-1)\) 上分子、分母同为负数，在 \((-1,2)\) 上异号，在 \((2,+\infty)\) 上同为正数. 因此不等式的解集为 \(\{x\mid x<-1\text{ 或 }x>2\}\).
\end{solution}

\begin{problem}
函数 \(y = 2\sin x\cos x - 1\)，\(x \in \R\) 的值域是\fillinblank{}.
\end{problem}

\begin{answer}
\([-2,0]\)
\end{answer}

\begin{solution}
利用 \(2\sin x\cos x=\sin2x\)，得 \(y=\sin2x-1\). 因为 \(-1\leqslant\sin2x\leqslant1\)，所以 \(-2\leqslant y\leqslant0\). 当 \(\sin2x=-1\) 或 \(1\) 时分别取得两个端点，故值域为 \([-2,0]\).
\end{solution}

\begin{problem}
双曲线 \( \frac{x^{2}}{m}-y^{2}=1 \) 上的点到左焦点的距离与到左准线的距离的比是 \(3\)，则 \(m\) 等于\fillinblank{}.
\end{problem}

\begin{answer}
\(\frac{1}{8}\)
\end{answer}

\begin{solution}
将双曲线写成标准形式 \(\frac{x^2}{a^2}-\frac{y^2}{b^2}=1\)，则 \(a^2=m\)，\(b^2=1\)，焦半距满足 \(c^2=a^2+b^2=m+1\). 双曲线的离心率为 \(e=\frac ca\)，而点到左焦点的距离与到左准线的距离之比等于 \(e\)，所以 \(e=3\). 因而 \(\frac{m+1}{m}=9\)，解得 \(m=\frac18\).
\end{solution}

\begin{problem}
如图，正四面体 \(ABCD\) 的棱长为 \(1\)，平面 \(\alpha\) 过棱 \(AB\)，且 \(CD \parallel \alpha\)，则正四面体上的所有点在平面 \(\alpha\) 内的射影构成的图形面积是\fillinblank{}.

\bitmapfigure[max width=0.55\linewidth]{img/2006/zhejiang_liberal/q14_fig1.png}
\end{problem}

\begin{answer}
\(\frac{1}{2}\)
\end{answer}

\begin{solution}
建立空间直角坐标系，取 \(A=(0,0,0)\)，\(B=(1,0,0)\)，\(C=\left(\frac12,\frac{\sqrt3}{2},0\right)\)，\(D=\left(\frac12,\frac{\sqrt3}{6},\sqrt{\frac23}\right)\). 令 \(\bs u=\overrightarrow{AB}=(1,0,0)\)，\(\bs v=\overrightarrow{CD}=\left(0,-\frac{\sqrt3}{3},\sqrt{\frac23}\right)\). 有 \(\lvert\bs u\rvert=\lvert\bs v\rvert=1\) 且 \(\bs u\perp\bs v\)，所以平面 \(\alpha\) 可取 \(\bs u\)，\(\bs v\) 为一组正交单位基底.

在这组基底下，四个顶点的正投影坐标分别为 \(A'=(0,0)\)，\(B'=(1,0)\)，\(C'=\left(\frac12,-\frac12\right)\)，\(D'=\left(\frac12,\frac12\right)\). 正四面体是四个顶点的凸包，正投影图形是四个投影点的凸包. 该图形的两条对角线 \(A'B'\)，\(C'D'\) 互相垂直，长度都为 \(1\)，所以面积为 \(\frac12\times1\times1=\frac12\).
\end{solution}

\section{解答题}
\begin{problem}
若 \(S_{n}\) 是公差不为 \(0\) 的等差数列 \(\{a_{n}\}\) 的前 \(n\) 项和，且 \(S_{1}\)，\(S_{2}\)，\(S_{4}\) 成等比数列.

\begin{enumerate}
\item 求数列 \(S_{1}\)，\(S_{2}\)，\(S_{4}\) 的公比；
\item 若 \(S_{2} = 4\)，求 \(\{a_{n}\}\) 的通项公式.
\end{enumerate}
\end{problem}

\begin{solution}
\begin{enumerate}
\item 设等差数列 \(\{a_n\}\) 的首项为 \(a_1\)，公差为 \(d\). 由等差数列前 \(n\) 项和公式，得 \(S_1=a_1\)，\(S_2=2a_1+d\)，\(S_4=4a_1+6d\). 因为 \(S_1\)，\(S_2\)，\(S_4\) 成等比数列，所以 \(S_2^2=S_1S_4\)，即 \((2a_1+d)^2=a_1(4a_1+6d)\). 整理得 \(d(d-2a_1)=0\). 由题设 \(d\neq0\)，所以 \(d=2a_1\)，从而 \(S_1=a_1\)，\(S_2=4a_1\)，\(S_4=16a_1\). 又因为 \(d\neq0\)，所以 \(a_1\neq0\)，数列 \(S_1\)，\(S_2\)，\(S_4\) 的公比为 \(4\).

\item 若 \(S_2=4\)，则 \(4a_1=4\)，得 \(a_1=1\)，且 \(d=2\). 因此 \(a_n=a_1+(n-1)d=1+2(n-1)=2n-1\).
\end{enumerate}
\end{solution}

\begin{problem}
如图，函数 \(y = 2\sin (\pi x + \varphi)\)，\(x \in \R\)，其中 \(0 \leqslant \varphi \leqslant \frac{\pi}{2}\) 的图象与 \(y\) 轴交于点 \((0,1)\).

\begin{enumerate}
\item 求 \(\varphi\) 的值；
\item 设 \(P\) 是图象上的最高点，\(M\)，\(N\) 是图象与 \(x\) 轴的交点，求 \(\overrightarrow{PM}\) 与 \(\overrightarrow{PN}\) 的夹角.
\end{enumerate}

\bitmapfigure[max width=0.50\linewidth]{img/2006/zhejiang_liberal/q16_fig1.png}
\end{problem}

\begin{solution}
\begin{enumerate}
\item 由图象与 \(y\) 轴交于点 \((0,1)\)，得 \(2\sin\varphi=1\). 又因为 \(0\leqslant\varphi\leqslant\frac{\pi}{2}\)，所以 \(\varphi=\frac{\pi}{6}\).

\item 函数取得最大值时，\(\pi x+\frac{\pi}{6}=\frac{\pi}{2}+2k\pi\)，所以图中最高点为 \(P\left(\frac13,2\right)\). 图中与 \(P\) 相邻的两个 \(x\) 轴交点满足 \(\pi x+\frac{\pi}{6}=k\pi\)，故 \(M\left(-\frac16,0\right)\)，\(N\left(\frac56,0\right)\).

于是 \(\overrightarrow{PM}=\left(-\frac12,-2\right)\)，\(\overrightarrow{PN}=\left(\frac12,-2\right)\). 设夹角为 \(\theta\)，则 \(\overrightarrow{PM}\cdot\overrightarrow{PN}=-\frac14+4=\frac{15}{4}\)，且 \(\lvert\overrightarrow{PM}\rvert=\lvert\overrightarrow{PN}\rvert=\frac{\sqrt{17}}2\). 因而 \(\cos\theta=\frac{15/4}{(\sqrt{17}/2)^2}=\frac{15}{17}\)，从而 \(\theta=\arccos\frac{15}{17}\).
\end{enumerate}
\end{solution}

\begin{problem}
如图，在四棱锥 \(P-ABCD\) 中，底面为直角梯形，\(AD \parallel BC\)，\(\angle BAD = 90^{\circ}\)，\(PA \perp\) 底面 \(ABCD\)，且 \(PA = AD = AB = 2BC\)，\(M\)，\(N\) 分别为 \(PC\)，\(PB\) 的中点.

\begin{enumerate}
\item 求证：\(PB \perp DM\)；
\item 求 \(BD\) 与平面 \(ADMN\) 所成的角.
\end{enumerate}

\bitmapfigure[max width=0.55\linewidth]{img/2006/zhejiang_liberal/q17_fig1.png}
\end{problem}

\begin{solution}
取长度单位使 \(PA=AD=AB=2\)，于是 \(BC=1\). 建立空间直角坐标系，使底面 \(ABCD\) 在 \(z=0\) 平面内，并取 \(A=(0,0,0)\)，\(B=(0,2,0)\)，\(D=(2,0,0)\)，\(C=(1,2,0)\)，\(P=(0,0,2)\). 于是 \(M=\left(\frac12,1,1\right)\)，\(N=(0,1,1)\).

\begin{enumerate}
\item 有 \(\overrightarrow{PB}=(0,2,-2)\)，\(\overrightarrow{DM}=\left(-\frac32,1,1\right)\)，因此 \(\overrightarrow{PB}\cdot\overrightarrow{DM}=0+2-2=0\)，所以 \(PB\perp DM\).

\item 由 \(\overrightarrow{AD}=(2,0,0)\)，\(\overrightarrow{AN}=(0,1,1)\)，可知平面 \(ADMN\) 的方程为 \(y=z\)，其一个法向量为 \(\bs n=(0,1,-1)\). 取 \(\overrightarrow{BD}=(2,-2,0)\)，设 \(BD\) 与平面 \(ADMN\) 所成的角为 \(\theta\)，则 \(\sin\theta=\frac{\lvert\overrightarrow{BD}\cdot\bs n\rvert}{\lvert\overrightarrow{BD}\rvert\lvert\bs n\rvert}=\frac{2}{2\sqrt2\cdot\sqrt2}=\frac12\). 由于线面角为锐角或直角，所以 \(\theta=\frac{\pi}{6}\).
\end{enumerate}
\end{solution}

\begin{problem}
甲，乙两袋装有大小相同的红球和白球，甲袋装有 \(2\) 个红球，\(2\) 个白球，乙袋装有 \(2\) 个红球，\(n\) 个白球. 现从甲，乙两袋中各任取 \(2\) 个球.

\begin{enumerate}
\item 若 \(n = 3\)，求取到的 \(4\) 个球全是红球的概率；
\item 若取到的 \(4\) 个球中至少有 \(2\) 个红球的概率为 \(\frac{3}{4}\)，求 \(n\).
\end{enumerate}
\end{problem}

\begin{solution}
\begin{enumerate}
\item 当 \(n=3\) 时，甲袋和乙袋中的球数分别为 \(4\) 和 \(5\). 从每袋任取 \(2\) 个球，所有取法等可能，取到的球全是红球的取法数为 \(\mathrm C_2^2\mathrm C_2^2\)，故所求概率为 \(\frac{\mathrm C_2^2\mathrm C_2^2}{\mathrm C_4^2\mathrm C_5^2}=\frac1{60}\).

\item 设从甲袋取出的红球数为 \(X\)，从乙袋取出的红球数为 \(Y\). 甲袋中有 \(P(X=0)=\frac{\mathrm C_2^0\mathrm C_2^2}{\mathrm C_4^2}=\frac16\)，\(P(X=1)=\frac{\mathrm C_2^1\mathrm C_2^1}{\mathrm C_4^2}=\frac23\). 乙袋中共有 \(n+2\) 个球，因此 \(P(Y=0)=\frac{\mathrm C_n^2}{\mathrm C_{n+2}^2}=\frac{n(n-1)}{(n+1)(n+2)}\)，\(P(Y=1)=\frac{\mathrm C_2^1\mathrm C_n^1}{\mathrm C_{n+2}^2}=\frac{4n}{(n+1)(n+2)}\).

至少有 \(2\) 个红球的对立事件是至多有 \(1\) 个红球. 由两袋取球相互独立，
\[
P(X+Y\leqslant1)=P(X=0)P(Y=0)+P(X=1)P(Y=0)+P(X=0)P(Y=1)
=\frac{n(5n-1)}{6(n+1)(n+2)}.
\]
由题意，\(P(X+Y\leqslant1)=1-\frac34=\frac14\)，所以 \(\frac{n(5n-1)}{6(n+1)(n+2)}=\frac14\)，即 \(7n^2-11n-6=0\). 解得 \(n=2\) 或 \(n=-\frac37\). 因为 \(n\) 是白球个数，故 \(n=2\).
\end{enumerate}
\end{solution}

\begin{problem}
如图，椭圆 \(\frac{x^{2}}{a^{2}} + \frac{y^{2}}{b^{2}} = 1\) （\(a > b > 0\)）与过点 \(A(2,0)\)，\(B(0,1)\) 的直线有且只有一个公共点 \(T\)，且椭圆的离心率 \(e = \frac{\sqrt{3}}{2}\).

\begin{enumerate}
\item 求椭圆方程；
\item 设 \(F_{1}\)，\(F_{2}\) 分别为椭圆的左，右焦点，求证：\(|AT|^{2} = \frac{1}{2} |AF_{1}| \cdot |AF_{2}|\).
\end{enumerate}

\bitmapfigure[max width=0.55\linewidth]{img/2006/zhejiang_liberal/q19_fig1.png}
\end{problem}

\begin{solution}
\begin{enumerate}
\item 由离心率条件 \(e=\frac ca=\frac{\sqrt3}{2}\)，得 \(c^2=a^2-b^2=\frac34a^2\)，所以 \(b^2=\frac14a^2\). 直线 \(AB\) 的方程为 \(\frac x2+y=1\)，即 \(y=1-\frac x2\). 将其代入椭圆方程，得 \(\frac{x^2}{a^2}+\frac{4(1-\frac x2)^2}{a^2}=1\)，即 \(2x^2-4x+4-a^2=0\). 直线与椭圆只有一个公共点，所以上述关于 \(x\) 的方程有两个相等的实根，从而 \((-4)^2-4\cdot2(4-a^2)=0\)，解得 \(a^2=2\)，进而 \(b^2=\frac12\). 所以椭圆方程为 \(\frac{x^2}{2}+2y^2=1\). 此时二次方程的重根为 \(x=1\)，代入直线方程得 \(T=(1,\frac12)\).

\item 椭圆的焦半距满足 \(c^2=a^2-b^2=\frac32\)，故 \(F_1=(-c,0)\)，\(F_2=(c,0)\). 因为 \(c<2\)，所以 \(\lvert AF_1\rvert=2+c\)，\(\lvert AF_2\rvert=2-c\)，从而 \(\lvert AF_1\rvert\lvert AF_2\rvert=(2+c)(2-c)=4-c^2=\frac52\). 另一方面，\(\lvert AT\rvert^2=(1-2)^2+\left(\frac12\right)^2=\frac54\). 因此 \(\lvert AT\rvert^2=\frac12\lvert AF_1\rvert\lvert AF_2\rvert\)，结论得证.
\end{enumerate}
\end{solution}

\begin{problem}
设 \( f(x) = 3ax^2 + 2bx + c \)，若 \( a + b + c = 0 \)，\( f(0)f(1) > 0 \).

\begin{enumerate}
\item 求证：方程 \( f(x)=0 \) 有实根；
\item 求证：\(-2 < \frac{b}{a} < -1\)；
\item 设 \(x_{1}\)，\(x_{2}\) 是方程 \(f(x) = 0\) 的两个实根，则 \(\frac{\sqrt{3}}{3} \leqslant |x_{1} - x_{2}| < \frac{2}{3}\).
\end{enumerate}
\end{problem}

\begin{solution}
由 \(a+b+c=0\)，得 \(c=-a-b\)，且 \(f(0)=c=-a-b\)，\(f(1)=3a+2b+c=2a+b\). 由 \(f(0)f(1)>0\)，得 \((-a-b)(2a+b)>0\)，即 \((a+b)(2a+b)<0\).

\begin{enumerate}
\item 若 \(a=0\)，则上式左边为 \(b^2\geqslant0\)，不可能小于 \(0\)，所以 \(a\neq0\). 令 \(r=\frac ba\). 方程 \(f(x)=0\) 的判别式为 \(\Delta=(2b)^2-4\cdot3a\cdot c=4\left(b^2+3a(a+b)\right)=4a^2(r^2+3r+3)\). 又 \(r^2+3r+3=\left(r+\frac32\right)^2+\frac34>0\)，所以 \(\Delta>0\)，方程有两个实根，因而有实根.

\item 因为 \(a\neq0\)，将 \((a+b)(2a+b)<0\) 两边同除以正数 \(a^2\)，得 \(\left(1+\frac ba\right)\left(2+\frac ba\right)<0\). 因此 \(-2<\frac ba<-1\).

\item 仍令 \(r=\frac ba\). 由根与判别式的关系，\(\lvert x_1-x_2\rvert^2=\frac{\Delta}{(3a)^2}=\frac49(r^2+3r+3)\). 由 \(-2<r<-1\)，有 \(r^2+3r+3=\left(r+\frac32\right)^2+\frac34\geqslant\frac34\)，从而 \(\lvert x_1-x_2\rvert^2\geqslant\frac13\)，即 \(\lvert x_1-x_2\rvert\geqslant\frac{\sqrt3}{3}\). 又因为 \((r+1)(r+2)<0\)，所以 \(r^2+3r+3<1\)，从而 \(\lvert x_1-x_2\rvert^2<\frac49\)，即 \(\lvert x_1-x_2\rvert<\frac23\). 综上，\(\frac{\sqrt3}{3}\leqslant\lvert x_1-x_2\rvert<\frac23\).
\end{enumerate}
\end{solution}
